\documentclass[12pt,a4paper,twoside,openright,titlepage,final]{article}
\usepackage{fontspec}
\usepackage{amsmath}
\usepackage{amsfonts}
\usepackage{amssymb}
\usepackage{makeidx}
\usepackage{graphicx}
\usepackage[hidelinks,unicode=true]{hyperref}
\usepackage[spanish,es-nodecimaldot,es-lcroman,es-tabla,es-noshorthands]{babel}
\usepackage[left=3cm,right=2cm, bottom=4cm]{geometry}
\usepackage{natbib}
\usepackage{microtype}
\usepackage{ifdraft}
\usepackage{verbatim}
\usepackage[nottoc]{tocbibind}
\usepackage{pdflscape}
\usepackage{fancyvrb}
\usepackage[obeyDraft]{todonotes}
\ifdraft{
	\usepackage{draftwatermark}
	\SetWatermarkText{BORRADOR}
	\SetWatermarkScale{0.7}
	\SetWatermarkColor{red}
}{}
\usepackage{booktabs}
\usepackage{longtable}
\usepackage{calc}
\usepackage{array}
\usepackage{caption}
\usepackage{subfigure}
\usepackage{footnote}
\usepackage{url}
\usepackage[titletoc]{appendix}

\setsansfont[Ligatures=TeX]{texgyreadventor}
\setmainfont[Ligatures=TeX]{texgyrepagella}
\setmonofont{FreeMono}

\usetikzlibrary{decorations.pathreplacing}

%*******************************************************
%                 NO MODIFICAR
\newcommand*{\FSfont}[1]{%
  \fontencoding{T1}\fontfamily{#1}\selectfont}

\newlength{\tpheight}\setlength{\tpheight}{0.9\textheight}
\newlength{\txtheight}\setlength{\txtheight}{0.9\tpheight}
\newlength{\tpwidth}\setlength{\tpwidth}{0.9\textwidth}
\newlength{\txtwidth}\setlength{\txtwidth}{0.9\tpwidth}
\newlength{\drop}
%*******************************************************

% Crea una portada con los siguientes parámetros
%
% #1 : Título 
% #2 : Subtítulo
% #3 : Subsubtítulo
% #4 : Autor(es)
% #5 : Lugar
%

\newcommand*{\portada}[5]{
\begin{titlepage}
\begingroup
\vspace*{1cm}
\drop = 0.2\txtheight
\centering
\vfill
{\Huge \scshape #1}\\[\baselineskip]
{\Large \textbf{#2}}\\[\baselineskip]
{\Large \scshape #3}\\[\baselineskip]
\vspace*{0.3cm}
{\large \textit{#4}}\\[0.5\drop]
\includegraphics[scale=0.35]{./imagenes/logoURJC.jpg}
\vspace*{1.5cm}

{\large \scshape #5, \today} \par
\begin{center}
\end{center}
\vfill\null
\endgroup
\end{titlepage}
}
 %*****************************************************
 


\author{José Ignacio Escribano}

\title{}
\setlength{\parindent}{0pt}

\begin{document}

\pagenumbering{alph}
\setcounter{page}{1}

\portada{Foro de preguntas}{Análisis de Big Data}{Preguntas Temas 8 y 9}{José Ignacio Escribano}{Móstoles}

\tableofcontents
\thispagestyle{empty}
\newpage

\pagenumbering{arabic}
\setcounter{page}{1}


\section{Preguntas}

\subsection{Hasta hace muy poco tiempo, la única opción para análisis de grafos escalable en Spark era la utilización de GraphX, cuyo catálogo completo de funciones y operaciones solo es accesible usando el lenguaje de programación Scala. Sin embargo, un nuevo elemento en Spark desde su versión 1.6 permite también acceder a estas funciones por primera vez con el lenguaje Python. Busque información al respecto e indique cómo se llama la abstracción que permite esta mejora y qué funciones ofrece.}

GraphFrames [1] es la abstracción que permite el manejo de grafos, de forma similar a cómo lo hace GraphX. GraphFrames se encuentra construido en la cima de Spark DataFrames, que hace que tenga las siguientes ventajas:

\begin{itemize}
	\item GraphFrames provee de APIs uniformes para Pyhton, Java y Scala. Por primera vez todos los algoritmos de GraphX están dsiponibles para Python y Java.
	\item GraphFrames permite usar las APIs de Spark SQL y DataFrames para hacer consultas.
	\item GraphFrames tiene total soporte de las mismas fuentes de datos que DataFrames, lo que permite leer y escribir grafos usando formatos como Parquet, JSON o CSV.
\end{itemize}

En GraphFrames, los vértices y aristas del grafo son representados como DataFrames, permitiendo almacenar información arbitraria tanto en vértices como en aristas.\\

Las principales funciones de la clase GraphFrame son:

\begin{itemize}
	\item bfs: búsqueda en amplitud del grafo.
	\item connectedComponents: calcula las componentes conexas del grafo.
	\item degrees: calcula el grado de cada vértice del grafo.
	\item pageRank: calcula el pageRank del grafo.
	\item shortestPaths: calcula los caminos más cortos de un conjunto de vértices del grafo.
	\item stronglyConnectedComponents: calcula las componentes fuertemente conexas del grafo.
	\item triangleCount: cuenta el número de triángulos del grafo.
\end{itemize}
Fuentes:\\
\begin{verbatim}
[1] http://graphframes.github.io/
[2] http://graphframes.github.io/api/python/graphframes.html
[3] https://databricks.com/blog/2016/03/03/introducing-graphframes.html
\end{verbatim}

\subsection{La API Spark SQL dentro del framework para big data Spark proporciona una abstracción de datos llamada DataFrames: http://spark.apache.org/docs/latest/sql-programming-guide.html ¿En qué otros lenguajes de programación piensa que se ha podido inspirar esta abstracción de datos? ¿Cuáles son las principales funciones que ofrecen los DataFrames?}

Un DataFrame es una colección de datos organizados por columnas con nombre. Son equivalentes a una tabla en una base de datos relacional o a un data frame de R o Pyhton, aunque con una mayor optimización que en estos dos lenguajes. Por lo que los DataFrames están inspirados en los homólogos de estos dos lenguajes.\\

Los DataFrames pueden ser creados a partir de distintas fuentes de datos como como archivos de datos estructurados, tablas en Apache Hive, bases de datos externas, o RDDs existentes.\\

La API de DataFrames se encuentra disponible para los lenguajes Scala, Java, Python y R.\\

Las principales características de los DataFrames son:

\begin{itemize}
	\item Escalado desde kilobytes hasta petabytes.
	\item Soporte para sistemas de almacenamiento y formatos de datos de array ancho.
	\item Optimización y generación de código a través del optimizador SparkSQL Catalyst. 
\end{itemize} 

Las principales funciones de la clase DataFrame (en Pyhton) son:

\begin{itemize}
	
	\item show: muestra el dataframe en pantalla. Ejemplo:
	
	\begin{verbatim}
	# Importamos la librería de Spark
	from pyspark.sql import SQLContext
	
	# df es un dataframe
	df.show()
	\end{verbatim}
	
	\item printSchema: imprime el esquema de un dataframe. Ejemplo:
	
	\begin{verbatim}
	df.printSchema()
	\end{verbatim}
	
	\item select: selecciona una columna
	
	\begin{verbatim}
		# Selecciona la columna name y la imoprime en pantalla
		df.select("name").show()
		
		# Selecciona la columna name, y a la variable age le suma 1.
		df.select(df["name"], df["age"]+1).show()
	\end{verbatim}
	
	\item filter: filtra información de un dataframe.
	
	\begin{verbatim}
		# Filtra los valores de la columna age que sean mayores de 21, y los muestra por pantalla.
		df.filter(df["age"] > 21)
	\end{verbatim}
	
	\item goupBy: agrupa por columna.
	
	\begin{verbatim}
		# Agrupa por la columna age, que cuenta y muestra por pantalla
		df.groupBy("age").count().show()
	\end{verbatim}
\end{itemize}

La lista completa se puede encontrar en [4].\\

La creación de un dataframe se consigue utilizando distintas funciones de la clase SQLContext. Por ejemplo, para leer desde un JSON, usamos el siguiente código.

\begin{verbatim}
# Creamos un SQLContext
sqlContext = SQLContext(sc)

# Crea el dataframe del archivo "people.json"
df = sqlContext.read.json("people.json")
\end{verbatim}

La lista completa de todos los formatos disponibles se pueden encontrar en [5].


Fuentes:\\
\begin{verbatim}
[1] http://spark.apache.org/docs/latest/sql-programming-guide.html
[2] https://databricks.com/blog/2015/02/17/introducing-dataframes-in-spark-for-large-scale-data-science.html
[3] https://databricks.com/blog/2015/06/02/statistical-and-mathematical-functions-with-dataframes-in-spark.html
[4] http://spark.apache.org/docs/latest/api/python/pyspark.sql.html#pyspark.sql.DataFrame
[5] http://spark.apache.org/docs/latest/api/python/pyspark.sql.html#pyspark.sql.SQLContext
\end{verbatim}

\end{document} 